\section{Threats to Validity}
\label{sec:analysis-threats}

The strongest validity support is internal consistency. The experiments use
fixed model weights, fixed input shape, functional parity validation,
repeated rounds, and a shared benchmark contract. The local and Kubernetes
stages also use CPU-affinity and single-threading controls where the runtime
permits. These choices follow systems benchmarking guidance that stresses
controlled comparisons, explicit measurement context, and caution when
interpreting performance numbers~\cite{benchmarking_sc15}.

The main limitation is external validity. The remaining threats below are
consolidated from the methodological scoping in \cref{ch:methods} and apply
to all stages of the staged pipeline.

\paragraph{Single-node clusters.}
The thesis-facing RQ1.4 \texttt{kind} and RQ1.5 AKS primary runs are both
single-node. Single-node deployments eliminate inter-node network hops and do
not capture the scheduling and networking variability of a multi-node
production cluster. RQ1.5b
(\cref{sec:rq15b}, \cref{sec:analysis-rq15b}) addresses this directly with a
multi-node AKS sensitivity stage that forces every chain hop across a node
boundary; the measured per-condition multi-node penalty is bounded between
+2.131~ms (\texttt{chain\_3svc}) and +4.086~ms (\texttt{chain\_5svc}), with
the monolithic baseline gaining +2.305~ms (\cref{tab:rq15b_delta}). The
single-node primary numbers should therefore be interpreted as a clean
architectural-overhead measurement that excludes the inter-node network
path; the multi-node sensitivity in \cref{sec:rq15b} provides a quantitative
bound on what production multi-node deployments add on top.

\paragraph{CPU-only inference.}
All stages use CPU inference with a single thread. GPU-accelerated inference
would substantially shift the compute distribution, reduce per-inference
time, and change the relative weight of the gRPC boundary cost. The results
are not directly generalisable to GPU-hosted microservice deployments.

\paragraph{Single model and workload.}
The thesis evaluates one model (ResNet-18) at batch size one with one input
shape. Different architectures, larger activation maps, or higher-batch
workloads may produce different cost curves and different optimal split
regions.

\paragraph{Controlled conditions versus production variability.}
The local stages apply strict CPU stabilisation and achieve low cross-round
variance. The Kubernetes and AKS stages are less controllable and exhibit
higher variability. Neither environment captures the full noise of a shared
multi-tenant cloud cluster, where CPU steal, network congestion, and noisy
neighbours can substantially affect tail latency.

\paragraph{Activation-protection motivation.}
He~et~al.\ \cite{he2021_attacking_collaborative} demonstrate that an untrusted
remote service holding only intermediate feature maps can reconstruct the
client's input, so in-transit encryption alone does not protect activations
at rest on the downstream segment. This motivates evaluating the
infrastructure-level protection of the downstream service (RQ2.2) in
addition to communication-level hardening (RQ2.1). The two layers address
different threat models: SEV-SNP defends against the cloud host, not against
an actively malicious application running inside the protected VM, so it is
not a direct solution to the malicious remote-service model in
He~et~al.

\paragraph{mTLS threat-model scope.}
Sidecar-based mTLS protects the inter-proxy communication path but does not
prevent plaintext from existing inside the endpoint after TLS termination.
The security claims for RQ2.1 are explicitly scoped to communication-level
hardening.

\paragraph{SEV-SNP threat-model scope.}
AMD SEV-SNP protects the VM boundary against the host platform but not
against a compromised guest OS or application process, and several
side-channel and availability threats fall outside its core guarantee. In
particular, recent work shows that confidential VMs remain vulnerable to
malicious-interrupt injection from the untrusted
hypervisor~\cite{schluter2024_heckler}; this thesis does not experimentally
evaluate such attacks. RQ2.2 therefore claims VM-level confidential
execution for the downstream segment (\texttt{service2\_only}), not
complete application-level isolation. Extending the confidential boundary
to both services in the chain is treated as future work.

\paragraph{Generalisation of the cited literature.}
Multi-node clusters, concurrent clients, GPUs, larger models, feature
compression, autoscaling, and different service-mesh configurations could
change the absolute overheads and possibly the preferred topology. The
thesis therefore supports bounded claims about the evaluated system rather
than a universal rule for all microservice inference deployments.
